%% NON MODIFICARE QUESTE IMPOSTAZIONI
\date{}
\usepackage[latin9]{inputenc}
\usepackage{graphicx}
\usepackage{fancyhdr}
\usepackage{ucs}
\usepackage{lastpage}
\usepackage{longtable}
\usepackage[colorlinks=true]{hyperref}
\hypersetup{
    colorlinks,
    citecolor=black,
    filecolor=black,
    linkcolor=black,
    urlcolor=black
}
\pagestyle{fancy}
\fancyhead{}
\fancyfoot{}
\newcommand{\HRule}{\rule{\linewidth}{0.5mm}}
\rhead{\doctitle \versiondoc - \TeamCommitted}
\lhead{\setlength{\unitlength}{1mm}
	\begin{picture}(0,0)
		\put(5,0){\includegraphics[scale=0.07]{img/LogoTeam}}
	\end{picture}}
\cfoot{\thepage\ di \pageref{LastPage}}

%%Comandi particolari per i nomi delle figure. Chiamare come una funzione LaTeX
\newcommand{\TeamCommitted}{\emph{Team Committed }}
\newcommand{\respProg}{\emph{Responsabile di Progetto }}
\newcommand{\ammProg}{\emph{Amministratore di Progetto }}
\newcommand{\analProg}{\emph{Analista }}
\newcommand{\verifProg}{\emph{Verificatore }}
\newcommand{\programProg}{\emph{Programmatore }}
\newcommand{\progetProg}{\emph{Progettista }}
\newcommand{\proponProg}{\emph{Proponente }}
\newcommand{\commitProg}{\emph{Committente }}
\newcommand{\propProg}{\emph{Dr. Amir Baldissera }}
\newcommand{\commProg}{\emph{Prof. Tullio Vardanega }}
\newcommand{\nameproject}{\emph{Progetto SafetyGame}}

%%DA QUI IN POI POTETE MODIFICARE

%% INSERIRE QUI IL NOME DEL DOCUMENTO (INSERITE SEMPRE UNO SPAZIO ALLA FINE DEL NOME)
\newcommand{\doctitle}{Norme di Progetto }

%% INSERIRE QUI LA VERSIONE ATTUALE DEL DOCUMENTO (INSERITE SEMPRE UNO SPAZIO ALLA FINE DELLA VERSIONE)
\newcommand{\versiondoc}{V2.0}

%%INSERITE QUI LA DATA DI COMPILAZIONE FINALE DEL DOCUMENTO
\newcommand{\datared}{2012/01/31}

%%INSERIRE QUI IL/I REDATTORI
\newcommand{\redattore}{\begin{itemize}
\item Giorgio Maggiolo
\item Alessandro Cornaglia
\end{itemize}}

%%INSERIRE IL/I NOME DEI VERIFICATORI CHE HANNO VERIFICATO IL DOCUMENTO
\newcommand{\verificatori}{\begin{itemize}
\item Lorenzo Braghetto
\item Massimo Dalla Piet�
\end{itemize}}

%%INSERIRE IL NOME DI CHI HA APPROVATO IL DOCUMENTO
\newcommand{\approvazione}{\begin{itemize}
\item Giorgio Maggiolo
\item Giacomo Quadrio
\end{itemize}}

%%INSERIRE LA TIPOLOGIA DI USO DEL DOCUMENTO [Interno/Esterno]
\newcommand{\usodoc}{Interno}

%%INSERIRE LA LISTA DI DISTRIBUZIONE DEL DOCUMENTO
\newcommand{\listadistr}{\TeamCommitted}

%%INSERIRE IL SOMMARIO DEL DOCUMENTO
\newcommand{\testosommario}{Questo documento vuol definire le norme volte alla regolamentazione del gruppo \TeamCommitted necessarie al processo di sviluppo del progetto SecurityGame}


